This section recalls measure-theoretic definitions, cites the standard \hyperref[measure.uniqueness-principle]{uniqueness theorem for Borel measures}, and proves the \hyperref[measure.translation-invariant-line]{real-line translation-invariance lemma}. The companion retains the checked interval and rectangle uniqueness results, together with more general supporting results that are no longer developed in the prose. Its coverage remains partial, and Steven's correspondence review is not complete.

The declaration \texttt{real\_borel\_measure\_uniqueness} matches the real-line case of the cited theorem: both measures are finite on compact sets and agree on every bounded half-open interval with distinct ordered endpoints. The declaration \texttt{rectangle\_borel\_measure\_uniqueness} matches its rectangle case on $\mathbb R\times(0,\infty)$, represented by a real coordinate and the subtype of positive real numbers. Both use Mathlib's genuine Borel measures; compact closures supply the required finite interval or rectangle masses. The manuscript refers readers to the standard uniqueness theorem rather than reproducing these proofs.

The shared module \texttt{TraceFormulas.MeasureTheory.BorelUniqueness} and the companion's \texttt{pi\_lambda}, \texttt{finite\_measure\_uniqueness}, and \texttt{measure\_uniqueness\_of\_finite\_cover} remain available as general supporting formalizations. They use Mathlib's Dynkin systems and generating families, including a countable finite-measure cover that need not be increasing. Mathlib's pi-system convention requires closure only for nonempty intersections; its Dynkin system contains the empty set and is closed under complements and countable disjoint unions. These technical results go beyond the shortened refresher. The interval proof specializes Mathlib's interval uniqueness theorem; the rectangle proof uses a countable cover by products of intervals and Borel generation in each coordinate.

The shared module \texttt{TraceFormulas.MeasureTheory.RealTranslation} uses Mathlib's Borel measure space on $\mathbb R$ and Lebesgue measure \texttt{volume}. Its assumptions are finiteness on compact sets and invariance of the measure under pushforward by every translation. Translations are measurable bijections, so this is the prose condition on translated Borel sets. The zero measure is included.

The declarations \texttt{translation\_invariant\_eq}, \texttt{translation\_invariant\_interval}, and \texttt{translation\_invariant\_unique\_scale} prove the measure equality, interval values including empty intervals, and uniqueness of a finite nonnegative scalar. The coefficient is the measure of $[0,1)$. The proof specializes Mathlib's general additive invariant-measure uniqueness theorem \texttt{measure\_add\_measure\_eq}; the prose instead subdivides intervals and uses rational approximation. This proof-route difference remains for Steven to review.

Mathlib's \texttt{Measure.Regular} includes finiteness on compact sets and compact inner approximation on open sets, while \texttt{Measure.InnerRegular} expresses compact inner approximation on all measurable sets. On $\mathbb R$, finiteness on compact sets supplies both via library instances. The companion does not yet formalize the refresher's complete set of examples, weighted-measure and pushforward identities, or the hyperbolic Radon measure.
